% Chapter 1
\chapter{مقدمه}

موضوع: طراحی و پیاده‌سازی چارچوب دوقلوهای دیجیتال چندعاملی - پیش‌بینی‌ و پایش بی‌درنگ داده‌های شبیه‌سازی موتورهای جت CMAPSS ناسا

مقدمه:
۱. بیان مسئله و ضرورت پژوهش (مشکلات ایجاد شده در صنایع بزرگ در هنگام خرابی دستگاه‌ها )
۲. اهمیت پژوهش (اهمیت پیش‌بینی زمان خرابی دستگاه‌های مختلف در صنعت و پیشگیری از آن)
۳. اهداف پژوهش (امکان پیاده‌سازی سامانه‌های پایش مبنی بر دوقلو‌های دیجیتال چندعاملی و اعمال تغییر با استفاده از مدل‌های زبانی بزرگ در صنایع مختلف)
۴. محدودیت‌های پژوهش و مجموعه داده CMAPSS (بیان مصنوعی بودن مجموعه داده اما بسیار استفاده شده در مقاله‌ها - بیان امکان پیاده‌سازی واقعی)
۵. خلاصه فصل‌ها

بررسی ادبیات و مبانی نظری:
مفاهیم پایه:
۱. مقدمه
۲. دوقلوی دیجیتال
۳. کاربرد‌های دوقلوی دیجیتال
۴. عامل
۵. سیستم‌های چندعاملی
۶. دوقلو‌های دیجیتال چندعاملی
۷. هوش مصنوعی در دوقلوهای دیجیتال
۸. یادگیری‌های عمیق و کاربرد آن در دوقلو‌های دیجیتال
۹. لایه توجه و سازوکار آن در مدل‌‌های یادگیری‌ عمیق
۱۰. شبکه‌های عصبی بازگشتی و کاربرد آن در دوقلو‌های دیجیتال (معرفی RNN و LSTM)
۱۱. مدل‌های زبانی بزرگ (معرفی ابزار‌های workflow در این زمینه مانند dify)
۱۲. جمع‌بندی
مرور کار‌های مرتبط:
۱. مقالات مرتبط با دوقلو‌های دیجیتال
۲. مقالات مرتبط با معماری دوقلو‌های دیجیتال
۳. مقالات مرتبط با دوقلو‌های دیجیتال چندعاملی
۴. مقالات مرتبط با کاربرد‌های مدل‌های زبانی بزرگ در دوقلو‌های دیجیتال
۵. مقالات مرتبط با داده ساختار CMAPSS

روش‌شناسی و معماری پیشنهادی
۱. رویکرد کلی (طراحی - پیاده‌سازی - ارزیابی)
۲. معماری پیشنهادی دوقلوی دیجیتال چندعاملی (جزئیات لایه‌ها به همراه شکل و دیاگرام جهت واضح‌تر شدن)
۳. توصیف مجموعه داده CMAPSS (ساختار - پیش‌پردازش‌ها)
۴. بررسی معماری پیشنهادی برای مجموعه داده CMAPSS (بررسی عامل‌های لایه کاربرد و کارکرد آن‌ها و توضیح انتراعی آن‌ها و ارتباط آن‌ها - ارتباط با مدل‌های زبانی - پایش بی‌درنگ - پیش‌بینی و ویژگی‌های آن - ...)
۵. جمع‌بندی

پیاده‌سازی و آزمایش
۱. مراحل پیاده‌سازی (در هر مرحله با چه ابزاری و گام به گام)
۲. عامل دوقلوی دیجیتال پایش آن(با استفاده از داده‌ CMAPSS - توضیح میکروسرویس مانیتورینگ موتور‌ها)
۳. عامل پیش‌بینی خرابی با استفاده از الگوریتم LSTM
۴. عامل مدل‌های زبانی بزرگ و تعامل با دوقلوی دیجیتال CMAPSS (جزییات پیاده‌سازی Dify و دستورات مورد قبول و مورد استفاده)
۵. عامل پایش بی‌درنگ دوقلوی مجازی و شبیه‌سازی ورود داده (جزئیات شبیه‌سازی بی‌درنگ و ارتباط آن با سایر سرویس‌ها)
۶. روش استقرار سرویس‌ها
۷. جمع‌بندی

ارزیابی و نتیجه‌گیری
۱. نتایج تجربی (نمودار‌ها و خروجی‌های LSTM و مقایسه نو‌آوری‌ها- هزینه‌های پیاده‌سازی)
۲. بحث و ارزیابی (نمونه‌های پرسش و پاسخ استفاده از مدل‌های زبانی برای دوقلوی دیجیتال - زمان پاسخ - نظرسنجی ... - ارزیابی نتیجه LSTM و مقایسه با سایر مقاله‌‌ها)
۳. محدودیت‌ها (پیاده‌سازی فقط وابسته به مجموعه داده - ...)
۴. نوآوری‌ها (معماری چند لایه‌ای - ادغام مدل‌های زبانی با دوقلو‌های دیجیتال با روش جدید و کم هزینه - در نظر گرفتن دلیل خروجی LSTM و ارائه مهم‌ترین سنسور‌های موثر در نتیجه عمر هر موتور- ترجیح پیش‌بینی کم‌تر به بیشتر)
۵. پیشنهادات برای کار‌های آینده
۶. پیشنهادات عملی
۷. جمع‌بندی


\section{مقدمه}

در حال حاضر، استفاده از دوقلوی دیجیتال\LTRfootnote{Digital Twins} در صنایع مختلف، از جمله تولید، بهداشت و درمان، حمل و نقل و حتی شهرسازی، در حال افزایش است. این تکنولوژی به سازمان‌ها این امکان را می‌دهد که فرآیندها و محصولات خود را به شکل مؤثری مدیریت کنند و تصمیمات هوشمندانه‌ای اتخاذ نمایند. به طور خاص، توانایی پیش‌بینی خرابی‌ها و مشکلات پیش‌رو از طریق تجزیه و تحلیل داده‌های جمع‌آوری شده، به شرکت‌ها، کارخانه‌ها و مدیران کمک می‌کند تا به سمت کاهش هزینه‌های عملیاتی و افزایش بهره‌وری حرکت کنند.


تاریخچه استفاده از شبیه‌ساز‌های دیجیتال به سال‌ها قبل برمی‌گردد. مشهورترین مثال از این کاربردها، سیستم شبیه‌سازی که در طی ماموریت آپولوی 13 توسط ناسا ایجاد شد، است. در این ماموریت، زمانی که یکی از مخازن اکسیژن در فضاپیما منفجر شد، تیم کنترل پرواز با استفاده از یک مدل شبیه‌سازی شده از فضاپیما و اجزای آن، توانستند راه‌حلی برای عبور از این مشکل پیدا کنند. این واقعه نشان داد که چطور شبیه‌سازی‌ها می‌توانند به مهندسان کمک کنند تا از راه دور روندهای پیچیده را در محیطی امن آزمایش کنند و چالش‌های زندگی واقعی را شبیه‌سازی کنند. هرچند این شبیه‌سازی در آن زمان، به‌عنوان دوقلوهای دیجیتال شناخته نمی‌شد، اما تأثیر آنها در درک عمیق‌تر از ارتباط میان فضای فیزیکی و دیجیتال غیرقابل انکار است.
\cite{b3}

در نهایت، با پیشرفت مستمر فناوری‌های دیجیتال، عصر جدیدی از نوآوری و کارایی در انتظار صنایع مختلف است. دوقلوی دیجیتال به عنوان یکی از مؤلفه‌های کلیدی این تحول، نقش مهمی در ارتباط و همگرایی میان دنیای فیزیکی و دیجیتال ایفا خواهد کرد. بدین ترتیب، می‌توان انتظار داشت که کاربردهای دوقلوی دیجیتال در آینده نزدیک گسترش یافته و به بهبود مستمر کیفیت خدمات و محصولات منجر شود.

در این پژوهش پس از اینکه به تعریف چنین مفاهیم مختلف در این زمینه پرداخته می‌شود، به کاربرد آنها در زمینه‌های مختلف اشاره شده و مثال‌هایی از هرکدام بیان می‌شود. سپس به بیان یک معماری پیشنهادی برای دوقلوی دیجیتال چندعاملی پرداخته شده و جزییات آن شرح داده می‌شود. در نهایت چند ابزار مناسب در استفاده از دوقلوی دیجیتال معرفی می‌گردد.




\section{دوقلوی دیجیتال چندعاملی}

در دوقلوی دیجیتال، ارتباط متقابل میان موجودیت دیجیتال و موجودیت فیزیکی و حقیقی، اهمیت بالایی دارد. از این رو باید دوقلوی دیجیتال را در واقع به این صورت تعریف کرد. به ایجاد یک همتای دیجیتال برای یک پدیده فیزیکی و حقیقی گفته می‌شود که نه‌ تنها تغییر در پدیده فیزیکی باعث تغییر در همتای دیجیتال آن می‌شود، بلکه تغییر در همتای دیجیتال نیز باعث تغییر در پدیده فیزیکی و حقیقی می‌شود. بر پایه این تعریف می‌توان یک معماری تازه برای دوقلوی‌های دیجیتال ارائه کرد که می‌توان آن را در چند لایه تعریف نمود.


لایه اول (لایه فیزیکی):

این لایه شامل اجزا و پدیده‌های فیزیکی است. هر شئ، جسم، پدیده، اتفاق، ابزار و چیزی که به صورت واقعی وجود دارد یا اتفاق می‌افتد و قرار است که دوقلوی دیجیتال آن‌ها ایجاد شود، در این لایه قرار می‌گیرند.


لایه دوم (لایه جمع‌آوری و خواندن اطلاعات):

این لایه شامل راه‌‌ و روش‌ها و ابزار‌هایی است که به کمک آن می‌توانیم اطلاعات و داده‌های مورد نیاز را از لایه قبل دریافت کرده و آن را آماده برای انتقال به لایه بعد می‌کند.

لایه سوم (لایه تبدیل):

این لایه از لایه قبل، داده‌های اولیه را دریافت کرده و آن را به اطلاعات معنا دار تبدیل می‌کند. برای این کار نیاز است تا داده‌های اولیه ابتدا تمیز گردند و سپس داده‌های مختلف با هم ترکیب شوند. این داده‌های تمیز شده و ادغام شده، آماده استفاده شدن هستند. (ETL)


لایه چهارم (لایه دوقلوی دیجیتال):

داده‌های آماده و تمیز شده وارد این لایه می‌شود و دوقلوی دیجیتال ایجاد می‌شود. در واقع هر موجودیتی که در لایه اول وجود داشت، در این لایه یک همتای دیجیتال پیدا می‌کند که به صورت مستقیم بر روی آن می‌تواند اثر گذارد. در واقع در این لایه می‌توان موجودیت‌ها و روابط بین آن‌ها را مشاهده کرد که قابل فهم توسط کاربر نیز است. یک روش نمایش برای این لایه می‌تواند استفاده از گراف باشد که به خوبی موجودیت و رابطه در آن قابل شرح و نمایش است.

لایه پنجم (لایه کاربرد/اپلیکیشن):

این لایه پس از لایه دوقلوی دیجیتال قرار دارد. این لایه می‌تواند عامل‌های مختلفی از نرم‌افزار‌ها و اپلیکیشن‌ها را در خود جای دهد که با استفاده از دوقلو‌های دیجیتال، کاربرد‌های مختلفی را ارائه دهد. کاربرد‌هایی مانند نمایش بی‌درنگ، سیمولیشن و شبیه‌سازی سه‌بعدی، آنالیز و تحلیل دوقلو‌ها، پیشگویی و پیش‌بینی، مانیتورینگ، تصمیم‌گیر، بررسی if scenario و ... . هر کدام از این عامل‌ها می‌توانند به صورت مستقل یا وابسته به هم کار کنند و به هدف خود برسند. برخی از این عامل‌ها مانند تصمیم‌گیری و پیش‌بینی می‌توانند با هم و به ترتیب اتفاق بیافتند و همچنین تاثیر خود را بر روی دوقلوی دیجیتال بگذارند. پس از اینکه تغییری در همتای دیجیتال رخ داد، چون مستقیما با لایه فیزیکی در ارتباط است، تاثیر بلافاصله بر روی همتای فیزیکی آن نیز اعمال می‌شود.








































